\begin{frame}{왜 이 구분이 중요한가}
  \begin{itemize}
    \item ``화면이 예쁘니까'' Isaac Sim을 고르면: 실제로는
      \textbf{물리 정밀도}가 필요한 작업을 \textbf{비싼 GPU}
      위에서 돌리는 셈이 되기 쉬움
    \item 반대로, 접촉력이 정확해야 하는 다리 로보를
      ``CartPole도 되는데''라며 \textbf{가장 단순한 환경}에서
      배우면: 13.1의 \kb{현실 격차(reality gap)}가 그 단순함에서
      \textbf{이미 박혀} 있다 --- 시뮬레이터 자체가 물리를 너무
      약하게 근사하고 있기 때문
  \end{itemize}
  \vskip0.6em
  \begin{columns}[T]
    \begin{column}{0.48\textwidth}
      \kb{축 A: 물리}
      \begin{itemize}
        \item 접촉력 정확도
        \item \textbf{로봇 제어}의 핵심
      \end{itemize}
    \end{column}
    \begin{column}{0.48\textwidth}
      \kb{축 B: 센서}
      \begin{itemize}
        \item 렌더링 리얼리즘
        \item \textbf{픽셀 입력 정책}의 핵심
      \end{itemize}
    \end{column}
  \end{columns}
\end{frame}
